% ===== tr-05 块2a：第 4 章 pp.39--43（原文页 39--43）=====
% [原文 p.39]
如果这个光滑化能被足够显式地描述，那么同样地，人们应当能够仅用组合的语言，提取出 $B$-模型计算所需的周期。

于是对数几何将在镜像对称的两侧都扮演重要角色；但正如上面所提示的，这两种计算的背后应当都有某些组合对象。

事实上，对数几何与热带几何密切相关。我们将在随后几节探讨热带几何如何同时控制上述图景的 A-模型侧与 $B$-模型侧，从而补全上面的图表：
\begin{center}
\begin{tikzpicture}[node distance=0.9cm]
\node[draw,ellipse,minimum width=2.5cm,minimum height=1.0cm] (A) {A-模型};
\node[draw,ellipse,right=of A,minimum width=3.1cm,minimum height=1.0cm] (L) {对数几何};
\node[draw,ellipse,right=of L,minimum width=2.5cm,minimum height=1.0cm] (B) {B-模型};
\node[draw,ellipse,below=1.4cm of L,minimum width=3.1cm,minimum height=1.0cm] (T) {热带几何};
\draw[<->] (A) -- (L);
\draw[<->] (L) -- (B);
\draw[<->] (L) -- (T);
\draw[<->] (T) -- (A);
\draw[<->] (T) -- (B);
\end{tikzpicture}
\end{center}
% [Tutor Note: 原文此图无图号、无图注，属示意图而非编号插图；此处按原样重排 TikZ，不写 \figtitle。图中五条箭头均为双向箭头。]

% [Tutor Note: $\operatorname{Trop}(0^\dagger)=\mathbb{R}_{\ge 0}$ 来自具体计算：$0^\dagger$ 的幽灵层是 $\mathbb{N}$，而 $\operatorname{Hom}(\mathbb{N},\mathbb{R}_{\ge 0})=\mathbb{R}_{\ge 0}$（把 $1\in\mathbb{N}$ 送到任一非负实数）。后文 p.40 中 $C_{\bar\eta}$ 的纤维计算都要用到这一点。]
\section{A-模型与热带几何}
对数几何与热带几何之间的第一个联系来自一个初等的组合构造。给定对数概形 $X^\dagger$，我们可以如下构造 $X^\dagger$ 的\emph{热带化}。对 $X$ 的每个几何点 $\bar\eta$，我们有一个幺半群 $\overline{\mathcal{M}}_{X,\bar\eta}$，从而有一个锥 $C_{\bar\eta}:=\operatorname{Hom}(\overline{\mathcal{M}}_{X,\bar\eta},\mathbb{R}_{\ge 0})$（这里 $\operatorname{Hom}$ 指幺半群同态，而 $\mathbb{R}_{\ge 0}$ 被赋予加法幺半群结构）。进一步，若 $\bar\eta$ 位于 $\bar\eta'$ 的闭包中，则有一条泛化映射 $\overline{\mathcal{M}}_{X,\bar\eta}\to\overline{\mathcal{M}}_{X,\bar\eta'}$。\footnote{由于必须在 étale 拓扑中工作，实际上可以有多条泛化映射。例如，若 $X$ 是结点三次曲线，则从结点 $\bar\eta$ 到一般点 $\bar\eta'$ 有两条泛化映射。}取对偶，这给出映射 $C_{\bar\eta'}\to C_{\bar\eta}$。若 $X$ 上的对数结构是精细的，则这些映射是严格凸有理多面体锥的面包含。于是我们可以通过作这些面包含所给出的等同，构成一个胞腔复形，从而得到一个多面体锥复形 $\operatorname{Trop}(X^\dagger)$。实际上，一般地这未必真的构成一个有意义的胞腔复形，因为泛化映射可能会在面上引出许多奇怪的自等同；但在我们此处想要描述的情形中，这不会造成问题。

这个构造是函子性的：若 $f:X^\dagger\to Y^\dagger$ 是对数概形的态射，则我们得到 $\operatorname{Trop}(f):\operatorname{Trop}(X^\dagger)\to\operatorname{Trop}(Y^\dagger)$。
% [原文 p.40]
例如，考虑环面退化 $\mathcal{X}\to D$ 的情形。如上一节所见，它给出对数概形的态射 $\mathcal{X}_0^\dagger\to 0^\dagger$。坏集 $Z\subseteq\mathcal{X}_0$ 正是 $\mathcal{X}_0$ 上对数结构不精细的轨迹。因此我们可以把上述热带化构造应用于 $\mathcal{X}_0^\dagger\setminus Z\to 0^\dagger$。现在 $\operatorname{Trop}(0^\dagger)=\mathbb{R}_{\ge 0}$ 是一条射线。另一方面，若 $x\in\mathcal{X}_0$ 是零维分层，则 $x$ 的局部邻域看起来像 $f_x:Y_x\to\mathbb{C}$ 在 $0$ 上的纤维，其中 $Y_x$ 是由某个格多胞形 $\sigma_x\subseteq M_{\mathbb{R}}$ 所定义的 $C(\sigma_x)\subseteq M_{\mathbb{R}}\oplus\mathbb{R}$ 给出的环面簇，而态射 $Y_x\to\mathbb{C}$ 由投影 $M_{\mathbb{R}}\oplus\mathbb{R}\to\mathbb{R}$ 给出。于是由习题 4.8.4 可知 $\overline{\mathcal{M}}_{\mathcal{X}_0,x}=C(\sigma_x)^\vee\cap(N\oplus\mathbb{Z})$，且 $C_x=C(\sigma_x)$。特别地，诱导映射 $\operatorname{Trop}(f):C_x\to\operatorname{Trop}(0^\dagger)$ 的纤维是 $\operatorname{Trop}(f)^{-1}(1)=\sigma_x$。由此容易验证
\[
\operatorname{Trop}(f):\operatorname{Trop}(\mathcal{X}_0^\dagger\setminus Z)\to\operatorname{Trop}(0^\dagger)=\mathbb{R}_{\ge 0}
\]
的纤维是
\[
\operatorname{Trop}(f)^{-1}(1)=B,
\]
其中 $B$ 带有多面体分解 $\mathcal{P}$。于是对偶相交复形来自一个非常一般的构造。特别是要注意，$B$ 只依赖于 $\mathcal{X}_0^\dagger$，不依赖于 $\mathcal{X}$（尽管不知道这一般构造时这也是显然的）。

现在让我们转向 A-模型；就本讨论而言，它指的是在 Calabi--Yau 流形上计数曲线。设我们有一个环面退化 $\mathcal{X}\to D$。我们想数一般纤维上的曲线。能否改为通过数 $\mathcal{X}_0$ 上的曲线来做到这一点，从而让问题具有更组合化的性质？

这个问题有很长的历史。关于这类问题的第一项工作属于 Li 与 Ruan [58] 以及 Ionel 与 Parker [46],[47]。他们本质上考虑的是如下情形：有一个退化 $\mathcal{X}\to D$，其特殊纤维 $\mathcal{X}_0=X_1\cup X_2$ 是两个光滑不可约分支的正规交叉并。他们证明了 $\mathcal{X}_0$ 有一套 Gromov--Witten 不变量理论，并且它给出的答案与一般纤维上的 Gromov--Witten 理论相同。此外，他们给出了粘合公式，断言 $\mathcal{X}_0$ 的 Gromov--Witten 不变量可以用两个对 $(X_i,X_1\cap X_2)$（$i=1,2$）的 Gromov--Witten 不变量来计算。这里，与对 $(X,D)$（其中 $D\subseteq X$ 是光滑除子）相伴的 Gromov--Witten 理论是\emph{相对} Gromov--Witten 不变量的理论，其中考虑 $X$ 中的曲线，并在曲线与 $D$ 的交点处施加若干给定的切触阶。这一粘合公式已被证明是 Gromov--Witten 理论中非常有力的工具。

2001 年，Bernd Siebert [74] 提出用对数几何来推广这些结果。与此同时，Jun Li 正在研究 Li--Ruan 与 Ionel--Parker 理论（它们是用辛技术实现的）的代数几何进路。他给出了相对 Gromov--Witten 不变量的一个令人满意的代数几何定义，并重新证明了粘合公式，其中用到了一些对数几何的技术。然而，这一理论存在一个技术困难。在 Gromov--Witten 理论中，按惯例允许
% [原文 p.41]
定义域曲线发展出泡（bubble）。但在相对 Gromov--Witten 理论中，还必须允许目标空间 $X$ 也发展出泡。当定义域曲线的一个不可约分支落入除子 $D$ 中时就会发生这种情况，此时与 $D$ 的切触阶就变得没有意义了。因此相对稳定映射的实际目标空间可能是 $X$ 在 $D\subseteq X$ 上粘接一串 $D$ 上的 $\mathbb{P}^1$-丛。这往往使分析更加困难，并且是把这些技术推广到更复杂的退化时的一个主要障碍。

这个问题的若干解决方案在 2011 年完成。Brett Parker 在 [66]、[67] 中提供了一个全新的范畴------\emph{爆裂流形}（exploded manifolds）范畴，用来研究 Gromov--Witten 理论。这些流形携带与对数空间类似的信息，但它是一个更灵活、「更软」的范畴，便于在其间工作。在 [67] 中他给出了这一设定下 Gromov--Witten 不变量的定义，并给出了粘合公式。此外，Siebert 与我 [34] 完成了对数 Gromov--Witten 不变量的一套理论；Abramovich 与 Chen [11],[1] 沿着 Siebert 最初的建议也做到了这一点。我在这里概述其基本想法。

\begin{definition}{4.9.1}
精细饱和对数概形 $W^\dagger$ 上的一个\emph{对数曲线}是一个精细饱和对数概形 $C^\dagger$，连同态射 $C^\dagger\to W^\dagger$，该态射相对维数为一、平坦、对数光滑，且所有几何纤维都约化。
\end{definition}

这里对数光滑性蕴含 $C\to W$ 的几何纤维是结点曲线，这很令人满意，因为这类曲线正是允许作为稳定映射定义域的那类曲线。对数结构也可以看作把标记点包含了进来。例如，给定 $W=\operatorname{Spec}\mathbb{C}$ 上的光滑曲线 $C$，可以取有限个点 $x_1,\dots,x_k\in C$，并给 $C$ 赋予与子集 $\{x_1,\dots,x_k\}\subseteq C$ 相伴的除子型对数结构。则 $C^\dagger$ 在 $W$ 上对数光滑，且 $W$ 上的对数结构是平凡对数结构 $\mathcal{M}_W=\mathcal{O}_W^\times$。

\begin{definition}{4.9.2}
设 $X^\dagger\to S^\dagger$ 是精细饱和对数概形的态射。$X^\dagger$ 中以 $W^\dagger$ 为底的一个\emph{对数曲线}是一条对数曲线 $C^\dagger/W^\dagger$，连同一个态射 $f:C^\dagger\to X^\dagger$，使得它们可嵌入对数概形的交换图表
\begin{center}
\begin{tikzcd}
C^\dagger \arrow[r,"f"] \arrow[d] & X^\dagger \arrow[d] \\
W^\dagger \arrow[r] & S^\dagger
\end{tikzcd}
\end{center}
\end{definition}

$X^\dagger$ 中的一条对数曲线称为\emph{稳定对数映射}，如果对每个几何点 $\bar w\to W$，$f$ 限制到底层带标记曲线 $C_{\bar w}\to\bar w$ 上都是一个通常的稳定映射。我们把数据记为 $(C^\dagger/W^\dagger,f)$。

这一定义可以按通常方式，通过给标记点加标号来进一步装饰。
% [Tutor Note: 这里 $\mathbb{N}^r$ 上的处理与 p.36 中 $\alpha_0:\mathbb{C}^\times\oplus\mathbb{N}\to\mathbb{C}$ 是同一机制：结构映射在 $\mathbb{N}^r$ 的非零元上取 $0$，故 $\mathcal{M}_C\oplus\mathbb{N}^r$ 仍是幺半群层而 $f^\#$ 仍由旧的 $f^\#$ 与包含合成。]
% [原文 p.42]
[34] 的主要工作是构造一个性态良好的稳定对数映射模空间。每当人们试图构造对数对象的模空间时，就会出现一个技术问题；Martin Olsson 在他的学位论文 [65] 中探讨过这个问题。问题如下。设给定一个稳定对数映射，其定义域为 $\pi:(C,\mathcal{M}_C)\to(W,\mathcal{M}_W)$。则 $\pi':(C,\mathcal{M}_C\oplus\mathbb{N}^r)\to(W,\mathcal{M}_W\oplus\mathbb{N}^r)$ 同样给出一个稳定对数映射的定义域。这里结构映射 $\alpha_C$（或 $\alpha_W$）在常值层 $\mathbb{N}^r$ 的非零元素上取值 $0$，而 $\pi'$ 在幺半群上的作用只是把 $\mathbb{N}^r$ 同构地映到 $\mathbb{N}^r$。新的映射 $f^\#$ 是旧的 $f^\#:f^{-1}\mathcal{M}_X\to\mathcal{M}_C$ 与包含 $\mathcal{M}_C\to\mathcal{M}_C\oplus\mathbb{N}^r$ 的复合。结果是，单个稳定对数映射会产生可数多个其他映射，因此稳定对数映射的叠根本不可能有限型，从而不可能紧合。

解决办法是对 $W$ 上在适当意义下通用的那些对数结构作等同。在上面的例子中，所涉及的所有对数曲线都来自如下对数概形的笛卡尔图表：
\begin{center}
\begin{tikzcd}
(C,\mathcal{M}_C\oplus\mathbb{N}^r) \arrow[r] \arrow[d] & (C,\mathcal{M}_C) \arrow[d] \\
(W,\mathcal{M}_W\oplus\mathbb{N}^r) \arrow[r] & (W,\mathcal{M}_W)
\end{tikzcd}
\end{center}
于是所有这些多余的对数曲线都可以看作由最初选定的对数曲线经对数基变换拉回而得到。

为解决这个问题，我们引入稳定对数映射的一个性质，称为\emph{基本}（basic）。我不想在此给出定义，因为它非常繁复；但基本稳定对数映射的重要性质是普适性与有界性，这由以下两个定理表达，它们是对 [34] 主要结果的总结：

\begin{theorem}{4.9.3}
给定稳定对数映射 $(C^\dagger/W^\dagger,f)$，存在基本稳定对数映射 $(C_b^\dagger/W_b^\dagger,f_b)$，可嵌入交换图表
\begin{center}
\begin{tikzcd}
C^\dagger \arrow[r] \arrow[d] & C_b^\dagger \arrow[r] \arrow[d] & X^\dagger \arrow[d] \\
W^\dagger \arrow[r] & W_b^\dagger \arrow[r] & S^\dagger
\end{tikzcd}
\end{center}
其中左方方块在精细饱和对数概形范畴中是笛卡尔的，且底层概形的映射 $W\to W_b$ 与 $C\to C_b$ 都是同构。此外，$(C_b^\dagger/W_b^\dagger,f_b)$ 以及上图中的诸映射由 $(C^\dagger/W^\dagger,f)$ 在相差唯一同构的意义下唯一确定。
\end{theorem}

\begin{theorem}{4.9.4}
设 $\mathcal{M}(X^\dagger/S^\dagger)$ 表示 $X^\dagger$ 中在 $S^\dagger$ 上的基本稳定对数映射的叠。则：
\begin{enumerate}
\item $\mathcal{M}(X^\dagger/S^\dagger)$ 是 Deligne--Mumford 叠。
% [Tutor Note: 图中 $W^\dagger=(\operatorname{Spec}\mathbb{C},Q\oplus\mathbb{C}^\times)$ 且 $\alpha(p,s)=s\ (p=0)$、$0\ (p\neq0)$；这与 p.36 中 $\alpha_0$ 的构造完全相同，而取 $Q=\mathbb{N}$（即 $\sigma_Q$ 为射线）时 $W^\dagger$ 就退化为标准对数点 $0^\dagger$，此时 $\sigma_Q^\vee=\mathbb{R}_{\ge 0}=\operatorname{Trop}(0^\dagger)$。]
% [原文 p.43]
\item 设 $\beta$ 表示一组选择：亏格 $g$、标记点个数 $k$、$H_2(X,\mathbb{Z})$ 中的一个同调类，以及标记点的一组切触数据（这一概念可以精确化）。令 $\mathcal{M}(X^\dagger/S^\dagger,\beta)$ 表示 $\mathcal{M}(X^\dagger/S^\dagger)$ 中由亏格 $g$、$k$ 个标记点的曲线构成、代表给定同调类并满足给定切触条件的基本稳定对数映射所构成的子叠。则在对 $X^\dagger$ 作某些技术性假设之下，若 $X$ 在 $S$ 上紧合，则 $\mathcal{M}(X^\dagger/S^\dagger,\beta)$ 在 $S$ 上紧合。
\item 进一步假设 $X^\dagger\to S^\dagger$ 对数光滑，则 $\mathcal{M}(X^\dagger/S^\dagger,\beta)$ 带有虚拟基本类，从而可以定义对数 Gromov--Witten 不变量。
\end{enumerate}
\end{theorem}

Abramovich 与 Chen 在 [1],[11] 中也得到了类似结果。

对于「完全在环面退化的中心纤维上工作以理解 A-模型」这一问题，这是一个有希望的开端。然而，这一理论中仍有两处重大空白需要填补。

第一，需要粘合公式的类比物。它应当使我们能把退化中心纤维上曲线的计算分解为更简单的片段。不过，这预计会相当微妙，目前仍是进行中的工作。我稍后会再多说一点人们期待这样的公式具有何种形式。

第二，如前所述，环面退化的中心纤维 $\mathcal{X}_0^\dagger\to 0^\dagger$ 只在集合 $Z$ 之外才是精细饱和的。因此上面关于稳定对数映射的定理全都不适用。很可能在这种情况下连稳定对数映射的定义都不是正确的定义。所以这一理论仍需推广。这也是 Michael Kasa 正在进行的工作。

让我们回到热带化函子。设有退化 $q:\mathcal{X}\to D$，我们假定它是对数光滑的（比如正规交叉退化），这样就不必担心 $\mathcal{X}$ 上对数结构不精细的奇异集 $Z$。照例，这给出 $\mathcal{X}_0^\dagger\to 0^\dagger$。设我们有一个点上的基本稳定对数映射，即有图表
\begin{center}
\begin{tikzcd}
C^\dagger \arrow[r,"f"] \arrow[d,"\pi"] & \mathcal{X}_0^\dagger \arrow[d,"q"] \\
W^\dagger=(\operatorname{Spec}\mathbb{C},Q\oplus\mathbb{C}^\times) \arrow[r,"g"] & 0^\dagger
\end{tikzcd}
\end{center}
这里 $Q$ 是由 $Q=\sigma_Q\cap\mathbb{Z}^n$ 给出的幺半群，其中 $\sigma_Q$ 是严格凸有理多面体锥；而 $W$ 上的对数结构由 $\alpha:Q\oplus\mathbb{C}^\times\to\mathbb{C}$ 给出，定义为
\[
\alpha(p,s)=\begin{cases} s & p=0,\\ 0 & p\neq 0.\end{cases}
\]
